%%%%%%%%%%%%%%%%%%%%%%%%%%%%%%%%%%%%%%%%%
% Islamic University of Technology (IUT) Assignment Title Page 
% LaTeX Template
% Version 1.0 (13/11/23)
%
% This template has been modified from:
% https://www.overleaf.com/latex/templates/uppsala-university-template/jvjprsfnzgbj
%
% Original author:
% WikiBooks (http://en.wikibooks.org/wiki/LaTeX/Title_Creation)
% Modified by Ishmam Tashdeed to fit IUT
% License:
% CC BY-NC-SA 3.0 (http://creativecommons.org/licenses/by-nc-sa/3.0/)

%----------------------------------------------------------------------------------------
%	PACKAGES AND OTHER DOCUMENT CONFIGURATIONS
%----------------------------------------------------------------------------------------

\documentclass[12pt]{article}
\usepackage[english]{babel}
\usepackage[utf8x]{inputenc}
\usepackage{amsmath}
\usepackage{lipsum}
\usepackage{graphicx}
\usepackage{float}
\usepackage[colorlinks = true,
linkcolor = blue,
urlcolor  = blue,
citecolor = blue,
anchorcolor = blue]{hyperref}
\usepackage{listings}
\setlength{\marginparwidth}{2cm}
\usepackage[colorinlistoftodos]{todonotes}
\usepackage{xcolor}

\definecolor{codegreen}{rgb}{0,0.6,0}
\definecolor{codegray}{rgb}{0.5,0.5,0.5}
\definecolor{codepurple}{rgb}{0.58,0,0.82}
\definecolor{backcolour}{rgb}{0.95,0.95,0.92}

\lstdefinestyle{mystyle}{
	backgroundcolor=\color{backcolour},   
	commentstyle=\color{codegreen},
	keywordstyle=\color{magenta},
	numberstyle=\tiny\color{codegray},
	stringstyle=\color{codepurple},
	basicstyle=\ttfamily\footnotesize,
	breakatwhitespace=false,         
	breaklines=true,                 
	captionpos=b,                    
	keepspaces=true,                 
	numbers=left,                    
	numbersep=5pt,                  
	showspaces=false,                
	showstringspaces=false,
	showtabs=false,                  
	tabsize=2
}

\lstset{style=mystyle}
\graphicspath{ {./images/} }

\begin{document}
	
	\begin{titlepage}
		
		\newcommand{\HRule}{\rule{\linewidth}{0.5mm}}
		
		\center
		
		%----------------------------------------------------------------------------------------
		%	HEADING SECTIONS
		%----------------------------------------------------------------------------------------
		
		\textsc{\LARGE Islamic University of Technology}\\[1.5cm]
		\includegraphics[scale=0.35]{iut_logo.png}\\[1cm]
		\textsc{\Large Server Programming}\\[0.5cm]
		\textsc{\large SWE 4538}\\[0.5cm]
		
		%----------------------------------------------------------------------------------------
		%	TITLE SECTION
		%----------------------------------------------------------------------------------------
		
		\HRule \\[0.4cm]
		{ \huge \bfseries Lab 1}\\[0.4cm]
		\large Building First REST API with FastAPI\\
		\HRule \\[1.5cm]
		
		%----------------------------------------------------------------------------------------
		%	AUTHOR SECTION
		%----------------------------------------------------------------------------------------
		
		\begin{minipage}{0.4\textwidth}
			\begin{centering} \large
				\emph{Author:}\\
				Mohammad Faiaj Jarif\\
				230042152\\
			\end{centering}
		\end{minipage}\\[2cm]
		
		%----------------------------------------------------------------------------------------
		%	DATE SECTION
		%----------------------------------------------------------------------------------------
		
		{\large \today}\\[2cm]
		
		\vfill
		
	\end{titlepage}
	
	\pagenumbering{roman}
	\tableofcontents
	\newpage
	\pagenumbering{arabic}
	\newpage
	
	%----------------------------------------------------------------------------------------
	%	INTRODUCTION
	%----------------------------------------------------------------------------------------
	\section{Introduction}
	
	In this lab, we built a small web server and explored how clients communicate with servers using HTTP. We created a simple API for \textbf{E-Shop}, an e-commerce application using \textbf{Python} and the \textbf{FastAPI} framework. The lab covered creating routes that handle GET and POST requests, working with path parameters, query parameters, and JSON data, and testing API endpoints using FastAPI's interactive documentation.
	
	\subsection{Learning Outcomes}
	
	By the end of this lab, we should be able to:
	
	\begin{enumerate}
		\item Explain the roles of a client, server, HTTP request, and HTTP response.
		\item Set up and run a FastAPI development server.
		\item Create routes that handle GET and POST requests.
		\item Work with path parameters, query parameters, and JSON data.
		\item Test API endpoints using FastAPI's interactive documentation.
		\item Explain why data stored in a Python list disappears when a server restarts.
	\end{enumerate}
	
	\newpage
	
	%----------------------------------------------------------------------------------------
	%	TASK A: GUIDED SETUP AND FIRST API
	%----------------------------------------------------------------------------------------
	\section{Part A: Guided Setup and First API}
	
	\subsection{Task A0 -- Check Python Installation}
	
	We verified that Python 3.10+ is installed and available from the terminal by running:
	
	\begin{lstlisting}[language=bash, caption=Checking Python version, numbers=none, frame=single, framesep=10pt, xleftmargin=1.5em, aboveskip=1.5em, abovecaptionskip=1em]
	python --version
	\end{lstlisting}
	
	\subsection{Task A1 -- Set Up Workspace}
	
	We created a folder named \texttt{eshop-lab1}, set up a virtual environment, and installed FastAPI:
	
	\begin{lstlisting}[language=bash, caption=Setting up the workspace, numbers=none, frame=single, framesep=10pt, xleftmargin=1.5em, aboveskip=1.5em, abovecaptionskip=1em]
		
	python -m venv .venv
	.venv\Scripts\Activate.ps1
	python -m pip install "fastapi[standard]"
		
	\end{lstlisting}
	
	\vfill
	
	\subsection{Task A2 -- Create First Server}
	
	We created \texttt{main.py} with a root endpoint that returns a welcome message.
	
	\begin{lstlisting}[language=Python, caption=Task A2 -- First server endpoint, numbers=none, frame=single, framesep=10pt, xleftmargin=1.5em, aboveskip=1.5em, abovecaptionskip=1em]
		
	from fastapi import FastAPI, HTTPException
		
	app = FastAPI()
		
	@app.get("/")
	def home():
		return {"message": "Welcome to Server Programming, Jarif"}
		
	\end{lstlisting}
	
	The server was started using:
	\begin{lstlisting}[language=bash, caption=Starting the development server, numbers=none, frame=single, framesep=10pt, xleftmargin=1.5em, aboveskip=1.5em, abovecaptionskip=1em]
		
	fastapi dev main.py
		
	\end{lstlisting}
	
	\begin{figure}[H]
		\centering
		\includegraphics[width=\textwidth]{img-1}
		\caption{Task A2 -- Welcome message at the root endpoint.}
		\label{fig:task_a2}
	\end{figure}
	
	\bigskip
	\bigskip
	
	%----------------------------------------------------------------------------------------
	%	TASK B: ROUTES AND PARAMETERS
	%----------------------------------------------------------------------------------------
	\section{Part B: Routes and Parameters}
	
	\subsection{Task B1 -- Personalized Greeting}
	
	We created a GET endpoint \texttt{/hello/\{name\}} that accepts a name as a path parameter and returns a personalized greeting.
	
	\begin{lstlisting}[language=Python, caption=Task B1 -- Personalized greeting endpoint, numbers=none, frame=single, framesep=10pt, xleftmargin=1.5em, aboveskip=1.5em, abovecaptionskip=1em]
		
	@app.get("/hello/{name}")
	def get_hello(name: str):
		return {"message": f"Hello, {name}!"}
		
	\end{lstlisting}
	
	\begin{figure}[H]
		\centering
		\includegraphics[width=\textwidth]{img-2} 
		\caption{Task B1 -- Personalized greeting endpoint output.}
		\label{fig:task_b1}
	\end{figure}
	
	\subsection{Task B2 -- Square of a Number}
	
	We created a GET endpoint \texttt{/square/\{number\}} that accepts an integer path parameter and returns the number along with its square.
	
	\begin{lstlisting}[language=Python, caption=Task B2 -- Square endpoint, numbers=none, frame=single, framesep=10pt, xleftmargin=1.5em, aboveskip=1.5em, abovecaptionskip=1em]
		
	@app.get("/square/{number}")
	def get_square(number: int):
		return {"number": number, "square": number * number}
	\end{lstlisting}
	
	\begin{figure}[H]
		\centering

		\includegraphics[width=\textwidth]{img-3}
		\caption{Task B2 -- Square endpoint output for a sample input.}
		\label{fig:task_b2}
	\end{figure}
	
	\newpage
	
	%----------------------------------------------------------------------------------------
	%	TASK C: FIRST E-SHOP API
	%----------------------------------------------------------------------------------------
	\section{Part C: First E-Shop API}
	
	We defined a product catalog stored in memory as a Python list with at least 10 entries:
	
	\begin{lstlisting}[language=Python, caption=Product catalog data, numbers=none, frame=single, framesep=10pt, xleftmargin=1.5em, aboveskip=1.5em, abovecaptionskip=1em]
		
	products = [
		{"id": 1, "name": "Laptop", "price": 75000},
		{"id": 2, "name": "Mouse", "price": 800},
		{"id": 3, "name": "Keyboard", "price": 1500},
		{"id": 4, "name": "Monitor", "price": 12000},
		{"id": 5, "name": "Headphone", "price": 2500},
		{"id": 6, "name": "Smartphone", "price": 35000},
		{"id": 7, "name": "Tablet", "price": 22000},
		{"id": 8, "name": "Printer", "price": 9500},
		{"id": 9, "name": "Webcam", "price": 1800},
		{"id": 10, "name": "USB Drive", "price": 700},
	]
	\end{lstlisting}
	
	\textbf{Note on code ordering:} In the actual \texttt{main.py} file, the endpoints for Tasks C1 and C2 are placed \emph{after} the endpoints for Tasks E1 and E2. This is because FastAPI matches routes in the order they are defined. If the parameterized route \texttt{/products/\allowbreak\{product\_id\}} were defined before \texttt{/products/\allowbreak search} and \texttt{/products/\allowbreak affordable}, it would capture the words ``search'' and ``affordable'' as product IDs and cause errors. By placing the fixed paths first, we avoid this routing conflict.
	
	
	\vfill
	
	\subsection{Task C1 -- List All Products}
	
	We implemented a GET endpoint \texttt{/products} that returns the entire products list.
	
	\begin{lstlisting}[language=Python, caption=Task C1 -- List all products, numbers=none, frame=single, framesep=10pt, xleftmargin=1.5em, aboveskip=1.5em, abovecaptionskip=1em]
		
	@app.get("/products")
	def get_products():
		return {"list": products}
	\end{lstlisting}
	
	\begin{figure}[H]
		\centering

		\includegraphics[width=25em, height=50em]{img-4} 
		\caption{Task C1 -- Listing all products.}
		\label{fig:task_c1}
	\end{figure}
	
	\subsection{Task C2 -- Get Product by ID}
	
	We implemented a GET endpoint \texttt{/products/\{product\_id\}} that returns a single product by its ID, or a 404 error if not found.
	
	\begin{lstlisting}[language=Python, caption=Task C2 -- Get product by ID, numbers=none, frame=single, framesep=10pt, xleftmargin=1.5em, aboveskip=1.5em, abovecaptionskip=1em]
		
	@app.get("/products/{product_id}")
	def get_product_by_id(product_id: int):
		
	for product in products:
		if product["id"] == product_id: 
			return product
		
	raise HTTPException(status_code=404, detail="Product not found")
	\end{lstlisting}
	
	\begin{figure}[H]
		\centering

		\includegraphics[width=\textwidth]{img-5}  % Replace example-image with task_c2
		\caption{Task C2 -- Fetching a product by its ID.}
		\label{fig:task_c2}
	\end{figure}
	
	\newpage
	
	%----------------------------------------------------------------------------------------
	%	TASK D: ADDING PRODUCTS
	%----------------------------------------------------------------------------------------
	\section{Part D: Adding Products with POST}
	
	\subsection{Task D1 -- Define the Request Body Model}
	
	We defined a Pydantic model \texttt{ProductCreate} to validate incoming JSON data for creating new products.
	
	\begin{lstlisting}[language=Python, caption=Task D1 -- Pydantic model for product creation, numbers=none, frame=single, framesep=10pt, xleftmargin=1.5em, aboveskip=1.5em, abovecaptionskip=1em]
	from pydantic import BaseModel
		
	class ProductCreate(BaseModel):
		name: str
		price: float
	\end{lstlisting}
	
	\subsection{Task D2 -- POST Endpoint to Add a Product}
	
	We implemented a POST endpoint \texttt{/addProduct} that accepts a JSON body, creates a new product with an auto-generated ID, appends it to the list, and returns the created product.
	
	\begin{lstlisting}[language=Python, caption=Task D2 -- Add product endpoint, numbers=none, frame=single, framesep=10pt, xleftmargin=1.5em, aboveskip=1.5em, abovecaptionskip=1em]
	@app.post("/addProduct")
	def insert_product(product: ProductCreate):
		
		new_product = {
			"id": len(products) + 1,
			"name": product.name,
			"price": product.price
		}
		
		products.append(new_product)
		return new_product
	\end{lstlisting}
	
	\begin{figure}[H]
		\centering

		\includegraphics[width=\textwidth]{img-6}
		\caption{Task D2 -- Adding a new product via POST request.}
		\label{fig:task_d2}
	\end{figure}
	
	\bigskip
	\bigskip
	
	%----------------------------------------------------------------------------------------
	%	TASK E: FILTERING AND SEARCH
	%----------------------------------------------------------------------------------------
	\section{Part E: Filtering and Search}
	
	\subsection{Task E1 -- Search Products by Name}
	
	We implemented a GET endpoint \texttt{/products/search} that accepts a query parameter \texttt{name} and returns products whose names contain the supplied text, ignoring capitalization.
	
	\begin{lstlisting}[language=Python, caption=Task E1 -- Search by name, numbers=none, frame=single, framesep=10pt, xleftmargin=1.5em, aboveskip=1.5em, abovecaptionskip=1em]
	@app.get("/products/search")
	def search_by_name(name: str):
		
		ans = []
		for product in products:
			if name.lower() in product["name"].lower():
				ans.append(product)
		
		return {"list": ans}
	\end{lstlisting}
	
	\begin{figure}[H]
		\centering

		\includegraphics[width=\textwidth]{img-7} 
		\caption{Task E1 -- Searching products by name.}
		\label{fig:task_e1}
	\end{figure}
	
	\bigskip
	\medskip
	\smallskip
	
	\subsection{Task E2 -- Filter Products by Price}
	
	We implemented a GET endpoint \texttt{/products/affordable} that accepts a query parameter \texttt{max\_price} and returns products with prices less than or equal to the given value.
	
	\begin{lstlisting}[language=Python, caption=Task E2 -- Filter by price, numbers=none, frame=single, framesep=10pt, xleftmargin=1.5em, aboveskip=1.5em, abovecaptionskip=1em]
	@app.get("/products/affordable")
	def filter_by_price(max_price: int):
		ans = []
		for product in products:
			if product["price"] <= max_price:
				ans.append(product)
		
		return {"list": ans}
	\end{lstlisting}
	
	\begin{figure}[H]
		\centering

		\includegraphics[width=\textwidth]{img-8}  
		\caption{Task E2 -- Filtering products by maximum price.}
		\label{fig:task_e2}
	\end{figure}
	
	\newpage
	
	%----------------------------------------------------------------------------------------
	%	CONCLUSION
	%----------------------------------------------------------------------------------------
	\section{Conclusion}
	
	In this lab, we successfully built a REST API for an E-Shop application using FastAPI. We learned how to:
	
	\begin{itemize}
		\item Set up a FastAPI development environment with a virtual environment.
		\item Create GET endpoints with path parameters and query parameters.
		\item Create a POST endpoint that accepts JSON request bodies validated by Pydantic models.
		\item Handle errors using \texttt{HTTPException} for cases like missing products.
		\item Manage route ordering to prevent parameterized routes from capturing fixed path segments.
	\end{itemize}
	
\end{document}